% GENERATED FILE. Edit canonical lessons, then run npm run generate:books.
% canonical-source-hash: fnv1a64:9850d9b19d9240a8
% canonical-lessons: KA-C259-gundi, KA-C259-svetar, KA-C259-salu2, KA-C259-karavastra, KA-C259-sara

\chapter{Button, Sweater, Shawl, Handkerchief, Necklace}
\label{ch:ka-button-sweater-shawl-handkerchief-necklace}
\hlchaptermodality{259}

\begin{chapteropening}
\textbf{By the end of this chapter:} \emph{I can say a button, a sweater, a shawl, a handkerchief and a necklace.}

Complete the last lesson of chapter 259: I can say a button, a sweater, a shawl, a handkerchief and a necklace.
\end{chapteropening}

\section[\texorpdfstring{\kn{ಗುಂಡಿ} (guṇḍi)}{ಗುಂಡಿ (guṇḍi)}]{\kn{ಗುಂಡಿ} (guṇḍi) — a button}

\label{lesson:KA-C259-gundi}



\begin{quote}
Before the new one: say the Kannada for a biscuit, then the Kannada for bread.
\end{quote}

\subsection*{You'll want to know: \kn{ಗುಂಡಿ}}
\textbf{\kn{ಗುಂಡಿ}} — \emph{guṇḍi} — \textquotedblleft{}a button\textquotedblright{}.

\begin{grammarlens}[title={What you've built}]
One more word for this chapter.
\end{grammarlens}

\subsection*{Your turn}
\begin{itemize}
\raggedright
  \item \emph{Say it:} \emph{guṇḍi}
  \item \emph{Say it:} \emph{guṇḍi}, once more
  \item \emph{Say it:} say \emph{guṇḍi}
  \item \emph{Recall:} say the Kannada for a biscuit, then the Kannada for bread, then say \emph{guṇḍi} again
\end{itemize}

\begin{tcolorbox}[breakable,colback=teal!4,colframe=teal!35!black,title={Before you move on}]
What does \kn{ಗುಂಡಿ} mean? (\textquotedblleft{}A button\textquotedblright{}.) Say it once more.
\end{tcolorbox}
\section[\texorpdfstring{\kn{ಸ್ವೆಟರ್} (sveṭar)}{ಸ್ವೆಟರ್ (sveṭar)}]{\kn{ಸ್ವೆಟರ್} (sveṭar) — a sweater}

\label{lesson:KA-C259-svetar}



\begin{quote}
Before the new one: say the Kannada for bread, then the Kannada for a button.
\end{quote}

\subsection*{You'll want to know: \kn{ಸ್ವೆಟರ್}}
\textbf{\kn{ಸ್ವೆಟರ್}} — \emph{sveṭar} — \textquotedblleft{}a sweater\textquotedblright{}.

\begin{grammarlens}[title={What you've built}]
One more word for this chapter.
\end{grammarlens}

\subsection*{Your turn}
\begin{itemize}
\raggedright
  \item \emph{Say it:} \emph{sveṭar}
  \item \emph{Say it:} \emph{sveṭar}, once more
  \item \emph{Say it:} say \emph{sveṭar}
  \item \emph{Recall:} say the Kannada for bread, then the Kannada for a button, then say \emph{sveṭar} again
\end{itemize}

\begin{tcolorbox}[breakable,colback=teal!4,colframe=teal!35!black,title={Before you move on}]
What does \kn{ಸ್ವೆಟರ್} mean? (\textquotedblleft{}A sweater\textquotedblright{}.) Say it once more.
\end{tcolorbox}
\section[\texorpdfstring{\kn{ಶಾಲು} (śālu)}{ಶಾಲು (śālu)}]{\kn{ಶಾಲು} (śālu) — a shawl}

\label{lesson:KA-C259-salu2}



\begin{quote}
Before the new one: say the Kannada for a button, then the Kannada for a sweater.
\end{quote}

\subsection*{You'll want to know: \kn{ಶಾಲು}}
\textbf{\kn{ಶಾಲು}} — \emph{śālu} — \textquotedblleft{}a shawl\textquotedblright{}.

\begin{grammarlens}[title={What you've built}]
One more word for this chapter.
\end{grammarlens}

\subsection*{Your turn}
\begin{itemize}
\raggedright
  \item \emph{Say it:} \emph{śālu}
  \item \emph{Say it:} \emph{śālu}, once more
  \item \emph{Say it:} say \emph{śālu}
  \item \emph{Recall:} say the Kannada for a button, then the Kannada for a sweater, then say \emph{śālu} again
\end{itemize}

\begin{tcolorbox}[breakable,colback=teal!4,colframe=teal!35!black,title={Before you move on}]
What does \kn{ಶಾಲು} mean? (\textquotedblleft{}A shawl\textquotedblright{}.) Say it once more.
\end{tcolorbox}
\section[\texorpdfstring{\kn{ಕರವಸ್ತ್ರ} (karavastra)}{ಕರವಸ್ತ್ರ (karavastra)}]{\kn{ಕರವಸ್ತ್ರ} (karavastra) — a handkerchief}

\label{lesson:KA-C259-karavastra}



\begin{quote}
Before the new one: say the Kannada for a sweater, then the Kannada for a shawl.
\end{quote}

\subsection*{You'll want to know: \kn{ಕರವಸ್ತ್ರ}}
\textbf{\kn{ಕರವಸ್ತ್ರ}} — \emph{karavastra} — \textquotedblleft{}a handkerchief\textquotedblright{}.

\begin{grammarlens}[title={What you've built}]
One more word for this chapter.
\end{grammarlens}

\subsection*{Your turn}
\begin{itemize}
\raggedright
  \item \emph{Say it:} \emph{karavastra}
  \item \emph{Say it:} \emph{karavastra}, once more
  \item \emph{Say it:} say \emph{karavastra}
  \item \emph{Recall:} say the Kannada for a sweater, then the Kannada for a shawl, then say \emph{karavastra} again
\end{itemize}

\begin{tcolorbox}[breakable,colback=teal!4,colframe=teal!35!black,title={Before you move on}]
What does \kn{ಕರವಸ್ತ್ರ} mean? (\textquotedblleft{}A handkerchief\textquotedblright{}.) Say it once more.
\end{tcolorbox}
\section[\texorpdfstring{\kn{ಸರ} (sara)}{ಸರ (sara)}]{\kn{ಸರ} (sara) — a necklace, a chain}

\label{lesson:KA-C259-sara}



\begin{quote}
Before the new one: say the Kannada for a shawl, then the Kannada for a handkerchief.
\end{quote}

\subsection*{You'll want to know: \kn{ಸರ}}
\textbf{\kn{ಸರ}} — \emph{sara} — \textquotedblleft{}a necklace, a chain\textquotedblright{}.

\begin{grammarlens}[title={What you've built}]
One more word for this chapter.
\end{grammarlens}

\subsection*{Your turn}
\begin{itemize}
\raggedright
  \item \emph{Say it:} \emph{sara}
  \item \emph{Say it:} \emph{sara}, once more
  \item \emph{Say it:} say \emph{sara}
  \item \emph{Recall:} say the Kannada for a shawl, then the Kannada for a handkerchief, then say \emph{sara} again
\end{itemize}

\begin{tcolorbox}[breakable,colback=teal!4,colframe=teal!35!black,title={Before you move on}]
What does \kn{ಸರ} mean? (\textquotedblleft{}A necklace, a chain\textquotedblright{}.) Say it once more.
\end{tcolorbox}
